\BourbakiExerciseGroup

¶ 1) 设 R 是紧空间 X 的开覆盖。证明存在 X 的一致结构的一致邻域 V，使对每个 \(x \in X\)，\(V(x)\) 包含在属于 R 的某个集中{[}注意：对每个 \(x \in X\) 存在一致邻域 \(W_x\) 使 \(\dot{W}_x(x)\) 包含在 R 的某个集中，并用有限多个集 \(W_x(x)\) 覆盖 X{]}。

\begin{enumerate}
\def\labelenumi{\arabic{enumi})}
\setcounter{enumi}{1}
\item
  证明拟紧空间 X 可一致化，当且仅当它的拓扑是紧空间 Y 的拓扑在一个满射 \(f: X \rightarrow Y\) 下的逆像。
\item
  设 \(f\) 是紧空间 \(\mathbf{X}\) 到紧空间 \(\mathbf{Y}\) 中的连续映射，\(\mathbf{V}\) 是 \(\mathbf{X}\) 的开一致邻域，使得关系 \(f(x) = f(y)\) 蕴含 \((x, y) \in \mathbf{V}\)。证明存在一致邻域 \(W\)（\(\mathbf{Y}\) 的），使关系 \((f(x), f(y)) \in W\) 蕴含 \((x, y) \in V\)。
\item
  设 X 是第一章 § 9 习题 II 中定义的局部紧空间。证明对角线 \(\Delta\) 在 \(X \times X\) 中的每个邻域都包含一个形如 \([x, b[ \times [x, b[\) 的集。由此推出 X 上只有唯一一个与 X 的拓扑相容的一致结构，且 X 关于这个一致结构的完备化可与区间 \(X' = [a, b]\)（赋有拓扑 \(\mathcal{C}v(X)\)）等同。（注意 X 上的每个超滤子关于任何与 X 的拓扑相容的一致结构都必是 Cauchy 滤子）。
\item
  证明一致空间 X 是准紧的，当且仅当 X 上的每个超滤子都是 Cauchy 滤子。（注意：若 X 不准紧，则存在 X 的一致邻域 V，使当 x 取遍 X 时诸集合 \(\mathbb{C}\mathrm{V}(x)\) 在 X 上生成一个滤子）。
\item
  设 X 是这样一个一致空间：其中每个点列都至少有一个聚点。证明 X 是准紧的（用反证法）。
\item
  子集 \(\mathbf{A}\)（一致空间 \(\mathbf{X}\) 的）称为有界的，如果对每个一致邻域 \(\mathbf{V}\)（\(\mathbf{X}\) 的），存在有限集 \(\mathbf{F}\) 和整数 \(n > 0\) 使 \(\mathbf{A} \subset \overset{n}{\mathbf{V}}(\mathbf{F})\)。
\end{enumerate}

\begin{enumerate}
\def\labelenumi{\alph{enumi})}
\item
  两个有界子集的并是有界的。有界子集的闭包是有界的。每个准紧子集都是有界的。
\item
  若 \(f: \mathbf{X} \to \mathbf{Y}\) 一致连续，则在 \(f\) 下，\(\mathbf{X}\) 的每个有界子集的像是 \(\mathbf{Y}\) 的有界子集。
\item
  在非空一致空间的积中，一个集是有界的当且仅当它的每个投影都是有界的。
\item
  对 X 的每个对称一致邻域 V，设 \(R_{V}\) 表示诸集合 \(\stackrel{n}{V}\)（当 n 取遍整数 \textgreater0）的并。证明 \(R_{V}\) 在 \(X \times X\) 中既开又闭，且用第 4 目的记号，
\end{enumerate}

\[
\mathbf {R} _ {\mathbf {v}} (x) = \mathbf {A} _ {x, \mathbf {v}}
\]

对每个 \(x \in \mathbf{X}\) 成立。假设 \(\mathbf{R}_{\mathbf{V}} = \mathbf{X} \times \mathbf{X}\) 对每个对称一致邻域 \(\mathbf{V}\)（\(\mathbf{X}\) 的）都成立（特别当 \(\mathbf{X}\) 连通时就是这种情形）；证明集 \(\mathbf{A} \subset \mathbf{X}\) 有界当且仅当对每个对称一致邻域 \(\mathbf{V}\)（\(\mathbf{X}\) 的）与每个 \(x_0 \in \mathbf{X}\)，存在整数 \(n > 0\) 使 \(\mathbf{A} \subset \mathbf{V}(x_0)\)。

\begin{enumerate}
\def\labelenumi{\arabic{enumi})}
\setcounter{enumi}{7}
\item
  设 X 是一致空间，V 是 X 的闭一致邻域，A 是 X 的紧子集。证明 V(A) 在 X 中是闭的（参看第一章 § 10 第 2 目定理 1 推论 5）。
\item
  设 X 是局部紧空间，它带有一个与其拓扑相容的一致结构，并且存在一致邻域 V 具有如下性质：\(V(x)\) 对所有 \(x \in X\) 都在 X 中相对紧。a) 证明 X 关于这个一致结构是完备的。
\end{enumerate}

\begin{enumerate}
\def\labelenumi{\alph{enumi})}
\setcounter{enumi}{1}
\item
  设 \(\mathbf{U}\) 是 \(\mathbf{X}\) 的对称一致邻域，使 \(\dot{\mathbf{U}}\subset\mathbf{V}\)。证明若 \(\mathbf{A}\) 在 \(\mathbf{X}\) 中相对紧，则 \(\mathbf{U}(\mathbf{A})\) 亦然。
\item
  证明 X 是仿紧的{[}利用 b)、习题 7 及第一章 § 9 第 10 目定理 5{]}。
\item
  设 W 是 X 的闭对称一致邻域，使 \(\dot{W} \subset V\)。证明若 A 在 X 中闭，则 \(\mathbf{W}(\mathbf{A})\) 亦然（利用习题 8）。
\end{enumerate}

\begin{enumerate}
\def\labelenumi{\arabic{enumi})}
\setcounter{enumi}{9}
\tightlist
\item
  设 X 是 Hausdorff 一致空间，它有一个一致邻域 V，使 \(\mathbf{V}(x)\) 对每个 \(x \in X\) 都是准紧的。证明完备化 \(\hat{X}\) 是局部紧的且满足习题 9 的条件。{[}若 W 是 X 的对称开一致邻域，使 \(\hat{W} \subset V\)，证明闭包 \(\overline{W}\)（W 在 \(\hat{X} \times \hat{X}\) 中的）满足：\(\overline{W}(x)\) 对每个 \(x \in \hat{X}\) 是紧的{]}。
\end{enumerate}

¶ 11) 设 X 是 Hausdorff 一致空间，U 是它的一致结构，F(X) 是 X 的所有非空闭子集组成的集。设 F(X) 赋有由 § 1 习题 5 a) 中定义的一致结构 U 诱导的一致结构。

\begin{enumerate}
\def\labelenumi{\alph{enumi})}
\item
  证明若 \(X'\) 准紧，则 \(\mathfrak{F}(\mathbf{X})\) 亦然。{[}若 \((\mathbf{A}_{i})\) 是 X 的由 V-小集（V 对称）组成的有限覆盖，证明诸集合 \(\tilde{\mathbf{V}}(\mathbf{B}_{k})\) 覆盖 \(\mathfrak{F}(\mathbf{X})\)，其中诸 \(B_{k}\) 取遍诸集 \(A_{i}]\) 的所有并。
\item
  点 \(\mathbf{A} \in \mathfrak{F}(\mathbf{X})\) 具有如下性质：\(\mathbf{A}\) 在 \(\mathfrak{G}(\tilde{\mathcal{U}})\) 于 \(\mathfrak{F}(\mathbf{X})\) 上诱导的拓扑中的每个邻域，都包含 \(\mathbf{A}\) 在 \(\mathfrak{G}_{\Theta}\)（第一章 § 8 习题 12）于 \(\mathfrak{F}(\mathbf{X})\) 上诱导的拓扑中的一个邻域，当且仅当 \(\mathbf{A}\) 准紧。
\item
  证明拓扑 \(\mathfrak{G}(\mathfrak{U})\) 与 \(\mathfrak{G}_{\Theta}\) 在集 \(\mathfrak{K}(\mathbf{X})\)（由 \(\mathbf{X}\) 的紧子集组成）上诱导相同的拓扑{[}利用 \(b\) ) 及 §1 习题 5{]}。
\end{enumerate}

\begin{enumerate}
\def\labelenumi{\arabic{enumi})}
\setcounter{enumi}{11}
\tightlist
\item
  \begin{enumerate}
  \def\labelenumii{\alph{enumii})}
  \tightlist
  \item
    设 R 是一个布尔代数，并把 R 与 R 的极大预滤子组成的集 \(\Omega\) 的子集所成的布尔代数等同于（第一章 § 6 习题 20）。对每个有限分拆 \(\overline{\omega} = (\mathbf{A}_i)\)（\(\Omega\) 的，由属于 R 的集组成），设 \(\mathbf{C}_{\overline{\omega}}\) 表示子集 \(\bigcup_{i} (\mathbf{A}_i \times \mathbf{A}_i)\)（\(\Omega \times \Omega\) 的）。
  \end{enumerate}
\end{enumerate}

证明诸集合 \(C_{\overline{\omega}}\) 构成 \(\Omega\) 上某个 Hausdorff 一致结构的一致邻域基本系。完备化 \(\hat{\Omega}\)（\(\Omega\) 关于这个一致结构的）是一个全不连通紧空间。\(\hat{\Omega}\) 的子集在 \(\hat{\Omega}\) 中既开又闭，当且仅当它形如 \(\overline{\mathbf{A}}\)，其中 \(\mathbf{A} \in \mathbf{R}\)。证明 \(\mathbf{A} \to \overline{\mathbf{A}}\) 是 \(\mathbf{R}\) 到 \(\hat{\Omega}\) 的既开又闭子集全体上的双射，并且 \(\overline{\mathbf{C}}\mathbf{A} = \overline{\mathbf{C}}\overline{\mathbf{A}}\)，\(\overline{\mathbf{A} \cup \mathbf{B}} = \overline{\mathbf{A}} \cup \overline{\mathbf{B}}\) 与 \(\overline{\mathbf{A} \cap \mathbf{B}} = \overline{\mathbf{A}} \cap \overline{\mathbf{B}}\)。由此推出典范映射 \(\Omega \to \hat{\Omega}\) 是双射。

\begin{enumerate}
\def\labelenumi{\alph{enumi})}
\setcounter{enumi}{1}
\tightlist
\item
  设 X 是一个 Hausdorff 拓扑空间，其中既开又闭的集构成拓扑的基。取 R 为由 X 的既开又闭子集组成的布尔代数；则 \(\hat{\mathbf{X}}\)（X 关于一致结构 \(\mathcal{U}\)（在 \(a\) ) 中定义的）的完备化）的拓扑在 X 上诱导的拓扑就是 X 上给定的拓扑。此外，X 上以开集组成的集为基的滤子全体的极大元素，以及 X 上以闭集组成的集为基的滤子全体的极大元素，都是关于一致结构 \(\mathcal{U}\) 的 Cauchy 滤子。由此推出存在连续满射 \(\varphi: \hat{\mathbf{X}} \to \hat{\mathbf{X}}, \psi: \mathbf{X}' \to \hat{\mathbf{X}}\)，其中 \(\hat{\mathbf{X}}, \mathbf{X}'\) 是第一章 § 9 习题 25 与 24 中定义的空间。*证明若 X 是有理直线 Q，则 \(\psi\) 不是双射。*若 X 是极不连通的（第一章 § 11 习题 21），则 \(\psi\) 是双射，\(\hat{\mathbf{X}}\) 是极不连通的，并且可与伴随 \(\mathbf{X}'\) 的半正则空间等同于（第一章 § 8 习题 20）；于是极不连通空间可以刻画为极不连通紧空间的稠密子空间。试由此给出一个没有孤立点的极不连通紧空间的例子{[}（参看第一章 § 11 习题 21 f){]}。
\end{enumerate}

特别地，若取 X 为离散空间，则 \(\hat{X}\) 可与 X 的超滤子空间等同于（第一章 § 9 习题 26）。

13) 设 X 是局部紧空间，U 是 X 的开子集。设 B 是 U 与 CU 的相对紧连通分支的并。证明 B 是开的，且是 U 与 CU 的那些在 CU 中既开又紧的子集的并。（把 X 嵌入它的单点紧化 \(X' = X \cup \{\omega\}\)，注意在 \(X' - U\) 中 \(\omega\) 的连通分支是 \(\{\omega\} \cup CB\)，并利用第 4 目命题 6）。

\begin{enumerate}
\def\labelenumi{\arabic{enumi})}
\setcounter{enumi}{13}
\item
  设 X 是紧空间，B 是 X 上由连通闭集组成的滤子基。证明 B 中诸集之（非空）交是闭且连通的（用反证法，利用第 3 目命题 4 及第一章 § 9 第 1 目定理 1）。
\item
  设 X 是紧空间。则 X 的非空闭子集组成的空间 \(\mathfrak{F}(\mathbf{X})\)，赋有由 \(\mathcal{C}(\tilde{\mathcal{U}})\) 诱导的拓扑（习题 11），是紧的（第一章 § 9 习题 13）。考察一个超滤子 \(\Psi\)（\(\mathfrak{F}(\mathbf{X})\) 上的），并假设对 X 的一致结构的每个一致邻域 V 与每个 \(\mathfrak{X} \in \Psi\)，存在 \(\mathbf{M} \in \mathfrak{X}\) 使 \(\mathbf{M}\) 的任何两点都可由一条包含在 \(\mathbf{M}\) 中的 V-链连接。证明 \(\Psi\) 在 \(\mathfrak{F}(\mathbf{X})\) 中的极限点 A 是 \(\mathbf{X}\) 的连通子集（利用第 4 目命题 6）。
\item
  设 X 是连通紧空间，A、B 是 X 的两个不相交的非空闭子集。证明存在 \(\mathbb{C}(\mathrm{A} \cup \mathrm{B})\) 的一个连通分支，其闭包与 A 和 B 都相交。（设 U 是 X 的一致结构的一致邻域，使 \(\overline{\mathbf{U}}(\mathbf{A}) \cap \overline{\mathbf{U}}(\mathbf{B}) = \varnothing\)。首先证明集 \(\mathfrak{M}_{\mathrm{U}}\)（即 \(\overline{\mathbb{C}(\mathrm{U}(\mathbf{A}) \cup \mathrm{U}(\mathbf{B}))}\) 的既与 \(\overline{\mathrm{U}(\mathbf{A})}\) 相交又与 \(\overline{\mathrm{U}(\overline{\mathrm{B}})}\) 相交的连通分支组成的集）非空，办法是证明：对每个对称一致邻域 W ⊂ U，存在 \(x \in \mathrm{U}(\mathbf{A})\) 与 \(y \in \mathrm{U}(\mathbf{B})\) 以及一条连接 x 与 y 的 W-链，其除 x 与 y 外的所有点都属于 \(\mathbb{C}(\mathrm{U}(\mathbf{A}) \cup \mathrm{U}(\mathbf{B}))\)，然后应用习题 15。对每个一致邻域 V ⊂ U，证明每个 K ∈ \(\mathfrak{M}_{\mathrm{V}}\) 包含一个集 H ∈ \(\mathfrak{M}_{\mathrm{U}}\)，并且集 \(\mathfrak{M}_{\mathrm{U},\mathrm{V}}\)（由诸集 H ∈ \(\mathfrak{M}_{\mathrm{U}}\)（包含在某个集 K ∈ \(\mathfrak{M}_{\mathrm{V}}\) 中的）组成）在 \(\mathfrak{F}(\overline{\mathrm{X}})\) 中是闭的；最后推出，当 V 取遍包含在 U 中的一致邻域时，诸集 \(\mathfrak{M}_{\mathrm{U},\mathrm{V}}\) 的交非空）。
\item
  设 X 是连通局部紧空间。设 K 是 X 的任一非空紧子集且 K ≠ X。证明 K 的每个连通分支都与 K 在 X 中的边界相交（利用第 4 目命题 6 的推论）。由此推出，若 A 是 X 的任一非空相对紧开子集且 A ≠ X，则 A 的每个连通分支的闭包与 CA 相交。
\end{enumerate}

¶ 18) a) 证明紧连通空间 X 不能是可数无穷多个互不相交的非空闭集的并{[}用反证法：若 (F \(_{n}\) ) 是 X 的一个由闭集组成的可数无穷分拆，借助习题 17 证明存在紧连通集 K，它不与 F \(_{1}\) 相交但与无穷多个集 F \(_{n}\) 相交；为此考察 F \(_{2}\) 的一点关于 F \(_{2}\) 的一个紧邻域（不与 F \(_{1}\) 相交的）的连通分支。然后使用归纳法{]}。

\begin{enumerate}
\def\labelenumi{\alph{enumi})}
\setcounter{enumi}{1}
\tightlist
\item
  在下列两个附加条件之一成立时，把 a) 的结果推广到局部紧连通空间 X：某个 \(F_{n}\) 是紧且连通的，或者 X 是局部连通的。{[}在第一种情形，借助习题 17 化归为 a){]}。
\end{enumerate}

\(^{*}c)\) 设 \(\mathbf{X}\) 是 \(\mathbb{R}^3\) 的子空间，它是下列子空间的并：\(A_n\) 是半直线 \(x > 0\)，\(y = 1/n\)，\(z = 0\)（ \(n \geqslant 1\)）；\(B_n\) 是区间 \(2n < x < 2n + 2\)，\(y = 0\)，\(z = 0\)（ \(n \geqslant 0\)）；\(C_n\) 是由关系式

\[
x = 2 n + 1, \quad 0 \leqslant y \leqslant \frac {1}{n + 1}, \quad z = y \left(y - \frac {1}{n + 1}\right),
\]

所定义的集，其中 \(n \geqslant 0\)。

证明 X 是连通且局部紧的，并且是可数无穷多个互不相交的连通非空闭集的并。*

¶ 19) 紧连通空间 X 称为在它的两点 x, y 之间不可约的，如果除 X 自身外，X 的紧连通子集中没有同时包含 x 与 y 的。

\begin{enumerate}
\def\labelenumi{\alph{enumi})}
\item
  若 \(x\) 与 \(y\) 是紧连通空间 \(X\) 的两个不同点，证明 \(X\) 有一个紧连通子空间 \(K\)（在 \(x\) 与 \(y\) 之间不可约的）（利用习题 14 与 Zorn 引理）。
\item
  证明至少有两点的紧连通空间不能在其每对点之间都不可约（利用习题 17）。
\item
  设 X 是紧连通空间，并设存在一点 \(a \in X\)，它有一个连通闭邻域 \(V \neq X\)。若 x 与 y 是 CV 的任意两点，证明存在 X 的紧连通子集，它包含 x 以及 a, y 两点之一而不含另一个（利用习题 17）。在同一假设下，证明若 x, y, z 是 X 的任意三个不同点，则 X 不能同时在 x 与 y 之间、y 与 z 之间以及 x 与 z 之间都不可约。
\end{enumerate}

\begin{enumerate}
\def\labelenumi{\arabic{enumi})}
\setcounter{enumi}{19}
\tightlist
\item
  设 X 是紧连通空间，并设 L 是 X 中具有连通邻域基本系的全体点组成的集。S = CL 的点称为 X 的奇异点。L 的内部 \(\dot{L}\) 的点以及 \(\bar{S}\) 的连通分支称为 X 的素成分。
\end{enumerate}

\begin{enumerate}
\def\labelenumi{\alph{enumi})}
\item
  设 \(\mathbf{A} \neq \mathbf{X}\) 是 \(\mathbf{X}\) 的非空开子集，设 \(\mathbf{K}\) 是 \(\overline{\mathbf{A}}\) 的一个连通分支，并设 \(\mathbf{F}\) 是 \(\overline{\mathbf{A}} \cap \mathbb{C}\mathbf{K}\) 的闭包。证明若 \(\mathbf{A} \cap \mathbf{K} \cap \mathbf{F}\) 非空，则其所有点都是奇异点。若 \(x \in \mathbf{A} \cap \mathbf{F} \cap \mathbf{K}\)，则连通分支 \(\mathbf{Q}\)（\(x\) 在 \(\mathbf{F}\) 中的）与 \(\mathbb{C}\mathbf{A}\) 相交{[}考察 \(\mathbb{C}\mathbf{K}\) 的一点（含于 \(V(x)\) 中的）及其在 \(\overline{\mathbf{A}}\) 中的连通分支，并利用习题 17 证明 \(\mathbf{F}\) 中可与 \(x\) 由 \(\mathbf{F}\) 中的一条 V-链连接的诸点所成的集与 \(\mathbb{C}\mathbf{A}]\) 相交。由此推出素成分 \(\mathbf{P}\)（包含 \(x\) 的）与 \(\mathbb{C}\mathbf{A}\) 相交（注意 \(\mathbf{P}\) 包含 \(x\) 在 \(Q \cap A\) 中的连通分支，并对紧连通空间 \(Q\) 及开子集 \(Q \cap A\)（\(Q\) 的）应用习题 17）。
\item
  由 a) 推出，若 P 是 X 的素成分，U 是 P 的邻域，则 P 有一个包含在 U 中的连通邻域（用反证法）。
\item
  设 V 是 X 的对称一致邻域，并用 \(V'\) 表示点对 \((x, y)\)（X 的点的）组成的集，这些点对可由一条 V-链连接，链中除 x 与 y 外的所有点都是奇异点。当 V 取遍 X 的对称一致邻域时，相应的诸集 \(V'\) 构成 X 上某个一致结构（一般不是 Hausdorff 的）的一致邻域基本系。设 \(X'\) 是伴随 Hausdorff 空间，并证明 \(X'\) 的诸点在典范映射 \(X \rightarrow X'\) 下的逆像是 X 的素成分。（紧连通）空间 \(X'\) 称为 X 的素成分空间。证明 \(X'\) 是局部连通的{[}利用 b) 与第一章 § 10 第 4 目命题 8{]}。
\end{enumerate}

\(^{*}d)\) 设 \((r_{n})\) 是含于 \(]0,1[\) 中的全体有理数按某个次序排成的序列。对每个无理数 \(x\in[0,1]\)，令

\[
f (x) = \sum_ {n} 2 ^ {- n} \sin 1 / (x - r _ {n}).
\]

设 X 是 f 的图像在 \(R^{2}\) 中的闭包。~证明 X 是连通的，且在其横坐标为 0 与 1 的两点之间不可约，但 X 只有一个素成分。*

*21) 在第一章 § 10 习题 20 b) 中定义的局部紧空间 X 中，设 S 是以 X 的连通分支为等价类的等价关系。证明商空间 X/S 不是 Hausdorff 的。*

\begin{enumerate}
\def\labelenumi{\arabic{enumi})}
\setcounter{enumi}{21}
\tightlist
\item
  证明每个全不连通紧空间 X 都同胚于有限离散空间的一个逆系统的逆极限。（考察 X 分成开集的有限分拆，并利用第 4 目命题 6 的推论）。
\end{enumerate}

¶ 23) 设 \(f: \mathbf{X} \to \mathbf{Y}\) 是连续开映射，其中 \(\mathbf{X}\) 局部紧，\(\mathbf{Y}\) 是 Hausdorff 的。

\begin{enumerate}
\def\labelenumi{\alph{enumi})}
\item
  若 K 是 Y 的任一紧连通子集，C 是 \(\overline{f}^{1}(K)\) 的任一紧连通分支，证明 \(f(C)=K\)（化归为 K=Y 的情形，并利用第 4 目命题 6 的推论）。*试给出一个非紧连通分支 \(C'\)（\(\overline{f}^{1}(K)\) 的）的例子，使得 \(f(C')\neq K_*\)
\item
  再假设 Y 局部紧且局部连通。证明若 U 是 Y 的任一连通开子集，R 是 \(\overline{f}^{1}(U)\) 在 X 中的任一相对紧连通分支，则 \(f(R)=U\)。{[}注意，在局部紧局部连通空间中，连通开集 U 是包含在 U 中且包含某定点 \(y_{0}\in U\) 的诸紧连通集的并；然后应用 a)。{]}
\item
  再假设 X 局部连通。证明若 K 是 Y 的任一连通子集（它是开集，或是内部非空的紧集），H 是 X 的任一紧子集，则 \(\overline{f}^{1}(K)\) 的包含在 H 中的连通分支只有有限多个。{[}利用 a) 与 b){]}。
\end{enumerate}

